% s09_value (abridged version). The value of the vertex.
\section{The value of the vertex}
\label{sec:val}

We compute the integral $\ZZ[g,0]=\Xiz_g(\emptyset)$ of \cref{def:set-fpot}, which is excluded from \cref{thm:van-main}. It is a pure Hodge integral: it depends only on the class $\Hg_{\bw}$ and the scale $E$. Recall $\beta_n(a,b)\coloneqq(n-1)!\,[\mathsf{t}^n]\,1/(4\sinh(a\mathsf{t}/2)\sinh(b\mathsf{t}/2))$.

\begin{thm}[{\Lg{Theorem~9.1}}]\genusstatus{AI-proposed, unchecked}\label{thm:val-main}
For $g\geq 2$,
\[
(-\epsilon_1\epsilon_2E^3)^{g-1}\,\ZZ[g,0]=\beta_{2g-2}(\epsilon_1,\epsilon_2).
\]
\end{thm}

In the normalisation of \cref{sec:van}, put $\mathsf{a}_1\coloneqq-\tau$, $\mathsf{a}_2\coloneqq1+\tau$, so that $\mathsf{a}_1\mathsf{a}_2=-\st$, and $\ZZ_h(\tau)\coloneqq\Xit_h(\emptyset)$.

\begin{thm}[{\Lg{Theorem~9.3}}]\genusstatus{AI-proposed, unchecked}\label{thm:val-tau}
For every $h\geq 2$, identically in $\tau$,
\[
\st^{\,h-1}\,\ZZ_h(\tau)=\beta_{2h-2}(\mathsf{a}_1,\mathsf{a}_2).
\]
\end{thm}

\Cref{thm:val-main} follows from \cref{thm:val-tau} by the scale law \eqref{eq:van-scale} at $r=0$, with $\tau=-\epsilon_1/E$, $\mathsf{a}_i=\epsilon_i/E$, $-\epsilon_1\epsilon_2E^3=\st E^5$ and the homogeneity $\beta_n(ca,cb)=c^n\beta_n(a,b)$. The proof of \cref{thm:val-tau} uses \cite[eq.~(1)--(3)]{LLZ03} only at $\mu=1^a$, through \cref{prop:van-real}; \cref{lem:van-dil}; the equation \eqref{eq:van-bil} at $k=2$; and textbook facts listed in \cref{sec:int-inputs}. It uses nothing from \cref{thm:van-main}.

\begin{proof}[Sketch of the proof of \cref{thm:val-tau}]
\emph{Step 1 (the slice).} Let $G_h(\yv)\coloneqq\sum_aM_h(1^a)\yv^a/a!$, with $M_h(1^a)$ as in \cref{sec:van-MV}. On the slice $p=(p_1,0,0,\dots)$, $p_1=\yv/(\st x)$, \eqref{eq:van-real} gives $\MVF=\sum_{h\geq 0}\st^{\,h-1}x^{2h-2}(G_h(\yv)-G_h(0))$.

\emph{Step 2 (a closed equation on the slice).} Frobenius' identity $\binom{n}{2}\chi_\nu(2,1^{n-2})=f_2(\nu)\chi_\nu(1^n)$ gives $p_1^2\de_2s_\nu=f_2(\nu)s_\nu$ on the slice, hence $\de_\tau\MVF=xp_1^2\de_2\MVF$ there: this is the cut-and-join eigenvalue. The equation \eqref{eq:van-bil} at $k=2$, divided by $\MVZ_0\MVZ_1$, reads $2\de_2(\MVF-\MVF_1)+\de_1^2(\MVF+\MVF_1)+(\de_1(\MVF-\MVF_1))^2=0$. Eliminating $\de_2$ gives, with $\thv=p_1\de_1$ and $\MVF_\tau\coloneqq\de_\tau|_p\MVF$,
\begin{equation}\label{eq:val-star}
\frac{2}{x}\bigl(1-\re^{\gamma_\tau\thv}\bigr)\MVF_\tau+(\thv^2-\thv)\bigl(1+\re^{\gamma_\tau\thv}\bigr)\MVF+\bigl((1-\re^{\gamma_\tau\thv})\thv\MVF\bigr)^2=0,
\end{equation}
an exact identity in $\bbQ(\tau)\fps{\yv}\lser{x}$ for the slice alone.

\emph{Step 3 (the dilaton point).} By \cref{lem:van-dil} at $\msU=\emptyset$, $G_h(\yv)=\pw^{2h-2}\mathsf{Q}_h(\pw-1)$ for $h\geq2$, a polynomial in $\pw$ with lowest coefficient $\mathsf{Q}_h(-1)=\ZZ_h$. In genus $0$ and $1$ the series $G_0$, $G_1$ are computed in closed form in $\pw$; genus $1$ uses $\int_{\Mbar_{1,1}}\psi=\int_{\Mbar_{1,1}}\lambda_1=1/24$. So \eqref{eq:val-star} holds coefficientwise in the monomials $x^a\pw^b$.

\emph{Step 4 (grading).} Under $x^a\pw^b\mapsto\check{x}^{\,b-a}N^{-b}$, with $\check{x}=1/x$, the operator $\re^{\gamma_\tau\thv}$ becomes the shift $f(N)\mapsto f(N+\tau)$ plus corrections of higher order in $\check{x}$, and $\MVF=x\mathfrak{A}+\mathfrak{B}+\check{x}\mathfrak{C}+\cdots$ with $\mathfrak{A}$ explicit and
\[
\mathfrak{B}=-\frac{N^2}{4\st}-\frac{1}{24}\log N+\mathfrak{B}_s,\qquad\mathfrak{B}_s\coloneqq\sum_{h\geq 2}\st^{\,h-1}\ZZ_hN^{2-2h}.
\]
The order $\check{x}^0$ of \eqref{eq:val-star} holds identically. The order $\check{x}^1$ does not involve $\mathfrak{C}$, and it is a linear inhomogeneous difference-differential equation $\LL[\mathfrak{B}_s]=(\text{explicit})$ for $\mathfrak{B}_s$ \Lg{eq.~(9.4)}, which involves only the $\ZZ_h$.

\emph{Step 5 (solution).} Under the bijection $\mathsf{t}^k\mapsto k!N^{-k-1}$, the shift becomes multiplication by $\re^{-\tau\mathsf{t}}$, and with $f_s\coloneqq\sum_{h\geq 2}\st^{\,h-1}\ZZ_h\mathsf{t}^{2h-2}/(2h-3)!$ the equation becomes $\mathfrak{L}(f_{\mathrm{sing}}+f_s)=0$, where $f_{\mathrm{sing}}\coloneqq-1/(\st\mathsf{t}^2)+(1+2\st)/(24\st)$ and, with $\valb\coloneqq1-\re^{-\tau\mathsf{t}}$,
\[
\mathfrak{L}f\coloneqq2\valb\Bigl(\frac{f_\tau}{\mathsf{t}}-\frac{f_{\mathsf{t}}}{\mathsf{a}_2}\Bigr)+(2-\valb)\frac{f}{\mathsf{a}_2}.
\]
The function $f_*\coloneqq1/(4\sinh(\mathsf{a}_1\mathsf{t}/2)\sinh(\mathsf{a}_2\mathsf{t}/2))$ satisfies $\mathfrak{L}f_*=0$, because $\valb\coth(\tau\mathsf{t}/2)=2-\valb$, and $f_*=f_{\mathrm{sing}}+f_{\mathrm{reg}}$ with $f_{\mathrm{reg}}=\sum_{h\geq 2}\beta_{2h-2}(\mathsf{a}_1,\mathsf{a}_2)\mathsf{t}^{2h-2}/(2h-3)!$. So $f_{\Delta}\coloneqq f_s-f_{\mathrm{reg}}$ satisfies $\mathfrak{L}f_{\Delta}=0$.

\emph{Uniqueness.} This equation alone does not determine the $\ZZ_h$: the general solution of $\mathfrak{L}f=0$ is $f_*\,\mathsf{g}((1+\tau)\mathsf{t})$ with $\mathsf{g}$ arbitrary. If $f_{\Delta}=b_n\mathsf{t}^n+\cdots$ with $b_n\neq 0$, the lowest order gives $b_n=C\tau^{-1}(1+\tau)^{n+1}$. But $f_{\Delta}$ is invariant under $\tau\mapsto-1-\tau$, which exchanges the weights $\tau$ and $-1-\tau$ of $\Hg_{\bwt}$, and $\mathsf{a}_1$ and $\mathsf{a}_2$; this forces $C=0$. So $f_{\Delta}=0$.
\end{proof}